\section{Structure \texttt{window\_config}}

La structure \texttt{window\_config} permet de configurer une fenêtre GTK4. Elle regroupe les paramètres géométriques, visuels et comportementaux nécessaires à sa création.

\subsection{Définition}

\begin{lstlisting}[language=C]
typedef struct {
    const char *title;
    const char *icon_name;
    const char *widget_name;
    const char *background_color;
    const char *background_image_path;
    int default_width;
    int default_height;
    int min_width;
    int min_height;
    int max_width;
    int max_height;
    bool resizable;
    bool decorated;
    bool modal;
    bool maximized;
    bool present_on_create;
    widget_style_config style;
} window_config;
\end{lstlisting}

\subsection{Description des champs}

\begin{center}
\begin{tabular}{|l|l|p{5cm}|}
\hline
\textbf{Champ} & \textbf{Type} & \textbf{Description} \\
\hline
title & const char* & Titre affiché dans la barre de titre \\
\hline
icon\_name & const char* & Icône du thème pour la barre de titre \\
\hline
widget\_name & const char* & Nom interne pour le débogage \\
\hline
background\_color & const char* & Couleur de fond (format CSS) \\
\hline
background\_image\_path & const char* & Chemin vers une image de fond \\
\hline
default\_width & int & Largeur par défaut en pixels \\
\hline
default\_height & int & Hauteur par défaut en pixels \\
\hline
min\_width & int & Largeur minimale \\
\hline
min\_height & int & Hauteur minimale \\
\hline
max\_width & int & Largeur maximale \\
\hline
max\_height & int & Hauteur maximale \\
\hline
resizable & bool & Autorise le redimensionnement \\
\hline
decorated & bool & Affiche la barre de titre \\
\hline
modal & bool & Fenêtre modale \\
\hline
maximized & bool & État maximisé à la création \\
\hline
present\_on\_create & bool & Affichage immédiat \\
\hline
style & widget\_style\_config & Configuration du style \\
\hline
\end{tabular}
\end{center}

\bigskip

Le champ \textbf{title} définit le texte affiché dans la barre de titre. Utilisé par les lecteurs d'écran pour l'accessibilité. La valeur peut être NULL.

Le champ \textbf{icon\_name} spécifie l'icône à afficher. Accepte les noms du thème système ("dialog-information", "document-open"). NULL utilise l'icône par défaut.

Le champ \textbf{widget\_name} définit un nom interne pour le débogage et les styles CSS. Apparaît dans les outils de développement GTK.

Le champ \textbf{background\_color} spécifie une couleur de fond au format CSS : noms ("red", "blue"), hexadécimal ("\#FF0000") ou RGB ("rgb(255, 0, 0)"). NULL utilise le thème par défaut.

Le champ \textbf{background\_image\_path} definit une image de fond. Le chemin doit pointer vers un fichier image valide. NULL desactive cette fonctionnalite.

Les champs \textbf{default\_width} et \textbf{default\_height} définissent les dimensions initiales en pixels.

Les champs \textbf{min\_width}, \textbf{min\_height}, \textbf{max\_width} et \textbf{max\_height} définissent les contraintes de dimensionnement. 0 désactive la limite.

Le champ \textbf{resizable} contrôle si l'utilisateur peut redimensionner la fenêtre. FALSE désactive les poignées de redimensionnement.

Le champ \textbf{decorated} contrôle la présence de la barre de titre et des bordures. FALSE crée une fenêtre sans barre de titre.

Le champ \textbf{modal} crée une fenêtre modale bloquant l'interaction avec les autres fenêtres jusqu'à sa fermeture.

Le champ \textbf{maximized} spécifie si la fenêtre est maximisée dès sa création.

Le champ \textbf{present\_on\_create} contrôle l'affichage immédiat après création.

Le champ \textbf{style} est une structure \texttt{widget\_style\_config} pour configurer le style.

\subsection{Exemple d'utilisation}

\begin{lstlisting}[language=C]
#include "window.h"

void create_main_window(GtkApplication *app) {
    window_config config = {
        .title = "Mon Application GTK",
        .default_width = 1024,
        .default_height = 768,
        .min_width = 800,
        .min_height = 600,
        .resizable = true,
        .decorated = true,
        .modal = false,
        .maximized = false,
        .present_on_create = true,
        .style = {
            .css_class = "main-window",
            .hexpand = true,
            .vexpand = true
        }
    };

    GtkWidget *window = create_window(app, &config);
    gtk_application_add_window(app, GTK_WINDOW(window));
}
\end{lstlisting}
